\section{Discussion}
\label{sec:discussion}

\para{Extending the Linear Representation Hypothesis.}
Our results demonstrate that the \lrh extends beyond clean binary concepts like truth and sentiment to culturally constructed categories with fuzzy boundaries. The concept ``is a Pok\'{e}mon'' is not a natural kind---it is defined by a commercial franchise with over 1,000 members spanning nine generations of games. Yet a simple difference-in-means direction captures this category with 88.4\% accuracy and large effect sizes (Cohen's $d$ up to 2.62). This suggests that linear representations are a general-purpose encoding mechanism in transformers, applicable even to arbitrary cultural categories learned from training data.

\para{Continuous rather than binary representations.}
A striking finding is that the model does not encode a sharp binary boundary between Pok\'{e}mon and non-Pok\'{e}mon. Instead, entities fall along a continuous spectrum, with the ranking matching human intuitions: game creatures from similar franchises are ``almost Pok\'{e}mon,'' while chairs and computers are ``maximally not Pok\'{e}mon.'' This graded representation is consistent with prototype theory in cognitive science \citep{rosch1975cognitive}, where category membership is a matter of degree rather than a sharp boundary.

\para{What drives Pok\'{e}mon quality?}
Three factors likely contribute to an entity's Pok\'{e}mon quality score. First, \textit{morphological similarity}: names with Pok\'{e}mon-like phonotactics (e.g., Crystalix, Thunderix) score high regardless of whether they refer to real entities. Second, \textit{semantic similarity}: creatures from similar game franchises co-occur with Pok\'{e}mon in training data, leading to similar representations. Third, \textit{franchise co-occurrence}: Samus, a Nintendo character who is not creature-like at all, scores high---likely because Nintendo franchise content frequently appears alongside Pok\'{e}mon content in the training corpus. Disentangling these factors is an important direction for future work.

\para{Limitations.}
Our study has several limitations. (1)~We analyze only GPT-2-XL; results may not generalize to other architectures or model scales. Larger models likely encode cleaner representations. (2)~We use a single prompt template (``\texttt{\{entity\} is a}''); other templates may activate different aspects of entity knowledge. (3)~We do not perform causal validation---we have not verified that adding the Pok\'{e}mon direction to non-Pok\'{e}mon activations produces Pok\'{e}mon-like model outputs. Without this, the direction may be correlated with but not causally responsible for Pok\'{e}mon-related behavior. (4)~The non-Pok\'{e}mon test set is small (55 examples), yielding wide confidence intervals. (5)~High Pok\'{e}mon quality scores for some entities may reflect training data co-occurrence rather than genuine conceptual similarity---the model may have seen these entities discussed alongside Pok\'{e}mon on gaming forums or wikis.

\para{Broader implications.}
Our findings have implications for creative AI applications. A model that encodes ``Pok\'{e}mon quality'' as a continuous, manipulable direction could, in principle, be steered to generate novel Pok\'{e}mon-like entities by adding the Pok\'{e}mon direction to the residual stream during generation \citep{turner2023activation}. More broadly, the ability to extract and rank entities along culturally defined axes opens possibilities for analyzing how models organize knowledge about brands, franchises, genres, and other human-constructed categories.